\documentclass[10pt]{article}
\usepackage[margin=1in]{geometry}
\usepackage{amsmath,amssymb}
\usepackage{booktabs}
\usepackage{microtype}
\usepackage[hidelinks]{hyperref}
\usepackage{titlesec}
\titlespacing*{\section}{0pt}{5pt}{2pt}
\titlespacing*{\subsection}{0pt}{3pt}{1pt}
\setlength{\parskip}{1.5pt}
\frenchspacing

\title{\vspace{-3.5em}Looking Monitorable Without Being Faithful:\\
Weak-to-Strong Reasoning Distillation Compresses Chain-of-Thought and Degrades Cue-Faithfulness}
\author{Abraham Yeung \\ \small CS 338 Project}
\date{\vspace{-1em}\small May 2026}

\begin{document}
\maketitle
\vspace{-1.5em}

\begin{abstract}
\noindent
Weak-to-Strong Reasoning (W2SR; Yuan et al.\ 2026) fine-tunes a capable student on
a \emph{weaker} model's chain-of-thought (CoT) and reports higher task accuracy.
Its safety case is incomplete: it measures only accuracy, never whether the
student's CoT stays \emph{monitorable}. We reproduce a W2SR capability gain and
extend it with a controlled CoT-monitorability study (Meek et al.\ 2025 cue-based
faithfulness eval, GPQA-Diamond). On a \emph{reasoning} student where faithfulness
has dynamic range, \textbf{W2SR significantly degrades cue-acknowledgment, $25\%\!
\to\!3\%$} (paired $\Delta\!=\!-0.22$, McNemar $p\!\approx\!2\times10^{-9}$), at
\emph{preserved} accuracy, while collapsing CoT length ${\sim}14\times$.
Decomposing the mechanism: the degradation is \textbf{compression-dominant}, but a
\textbf{real, moderate residual persists at matched length} (acknowledgment
$44\%\!\to\!19\%$, $p\!=\!0.016$; directionally echoed by three further
length-controlled cuts)---the distilled model surfaces decision-relevant cues less
\emph{even per unit reasoning}. A teacher-strength axis shows the effect is
\emph{compression-general}: a strong (14B) teacher degrades faithfulness nearly as
much as a weak (1.5B) one. On a non-reasoning instruct student faithfulness is
floored, so the effect is undetectable there---an artifact of substrate, which we
use to motivate the reasoning-student test. Reproduction itself required an
in-family weak teacher and is genuine only on a base student (a headroom probe
separates elicitation from prompting). \textbf{A capability-only certification of a
weak-to-strong student can pass a model whose CoT has become markedly less
revealing.}
\end{abstract}

\vspace{-0.5em}
\section{Introduction}
Chain-of-thought ``monitorability'' is the hope that reading a model's reasoning
reveals \emph{why} it answered, so an overseer can catch undesired influences.
Weak-to-Strong Reasoning (W2SR)~\cite{yuan2026} reports that fine-tuning a strong
student on a \emph{weaker} model's CoT \emph{raises accuracy}. Faithfulness
work~\cite{turpin2023,meek2025,chua2025} shows models often do not verbalize the
cues that actually move their answers.

\textbf{Critique.} W2SR certifies the student on accuracy alone and is silent on
whether the distilled CoT stays monitorable. If distillation changed \emph{how} a
model reasons on the page without changing \emph{what it reveals}, an accuracy
check would miss it.

\textbf{Contributions} (reproduce / critique / extend):
\textbf{(1)} we reproduce a W2SR capability gain on MATH and show it is conditional
---it requires an \emph{in-family} weak teacher, and is genuine only on a
\emph{base} student (on an instruct student it is a CoT-prompting artifact;
\S\ref{sec:repro}); \textbf{(2)} holding capability fixed, we measure CoT
monitorability and find that on a reasoning student W2SR \textbf{significantly
degrades cue-faithfulness}, predominantly by compressing the CoT but with a real
matched-length residual, and that this is \emph{compression-general} across teacher
strength (\S\ref{sec:ext}).

\section{Setup}
\textbf{W2SR}~\cite{yuan2026}: sample CoT from a weak teacher on a problem set, keep
correct + incorrect traces, LoRA-SFT a stronger student; report held-out Pass@1.
\textbf{Monitorability}~\cite{meek2025}: each question is paired with a \emph{cue}
(authority opinion, leaked answer key, embedded metadata, etc.) pointing at a
target; an LLM judge reads the full CoT and scores \emph{faithfulness}
(does the CoT \emph{acknowledge} the cue) and \emph{verbosity} (a
factor-utilization rate). We use the five cues shipped in the released eval and a
single judge (\texttt{claude-sonnet-4-6}, $T{=}0$); faithfulness is reported as
acknowledgment among \emph{cued} samples. All evaluation is greedy ($T{=}0$); trace
sets are hashed. Training is LoRA; generation/training/gating run on Modal GPUs.

\section{Reproduction: when does W2SR fire?}
\label{sec:repro}
With a foreign-style weak teacher (DeepSeek-R1-distilled 1.5B, or SimpleRL-Zoo)
distilled into a Qwen2.5-7B student, we found \emph{no} Pass@1 gain and frequent
degeneration, across students, ranks, and decodings. W2SR reproduces only with an
\textbf{in-family, native-Qwen} weak teacher: \texttt{Qwen2.5-Math-1.5B-Instruct}
($\sim$67\% on the train problems) $\to$ \texttt{Qwen2.5-Math-7B} lifts held-out
MATH (L3--5) Pass@1 from $0.325$ to $0.645$, format-valid $1.00$.

\textbf{Genuine only on a base student (headroom probe).} A $+5$pt gate-pass on an
already-CoT-capable model may be re-arrangement, not elicitation. Scoring the
\emph{untrained} model under a zero-shot CoT prompt (the prompting-only ceiling)
and reporting $\text{W2SR}-\text{(zero-shot CoT)}$ (Table~\ref{tab:repro}): on the
base student W2SR adds $+0.41$ \emph{beyond prompting} (genuine latent-capability
elicitation); on an instruct student it adds $+0.00$---its apparent gain over an
unelicited baseline is purely the CoT-elicitation confound. Two conditions Yuan et
al.\ leave implicit: in-family teacher, and a base (non-CoT-eliciting) student.

\begin{table}[t]\centering\small
\caption{Reproduction (held-out MATH L3--5 Pass@1). In-family native-Qwen teacher;
headroom probe separates genuine elicitation from CoT-prompting. Genuine on the
base student, artifact on the instruct student.}
\label{tab:repro}
\begin{tabular}{lcccc}
\toprule
student & unelicited & zero-shot CoT & W2SR & W2SR $-$ CoT \\
\midrule
Qwen2.5-Math-7B (base) & 0.325 & 0.240 & 0.645 & $\mathbf{+0.405}$ (genuine)\\
Qwen2.5-7B-Instruct & 0.230 & 0.630 & 0.630 & $\mathbf{+0.000}$ (artifact)\\
\midrule
\multicolumn{5}{l}{\footnotesize Foreign-style teachers (R1-distill-1.5B, SimpleRL-Zoo) $\to$ 7B: no gain / collapse.}\\
\bottomrule
\end{tabular}
\end{table}

\section{Extension: does W2SR preserve CoT monitorability?}
\label{sec:ext}
\textbf{Capability-controlled design.} We compare a baseline student, a W2SR student
(LoRA-SFT'd on a weak in-family teacher's MATH CoT), and references, on GPQA. The
SFT lands the student at its own CoT ceiling, so faithfulness differences are not a
capability confound; the question is ``does distillation change faithfulness,'' not
``does it improve capability.''

\textbf{Instruct student is floored (motivation).} On \texttt{Qwen2.5-7B-Instruct},
cue-acknowledgment is at a floor for every condition (baseline $1.8\%$ at $N{=}970$;
W2SR $3.6\%$; strong-teacher control $4.6\%$), while teacher references reach
$22$--$37\%$. We cannot detect faithfulness transfer here---but verbosity transfers
strongly (W2SR$-$baseline $\Delta\!=\!+0.26$, CI $[0.22,0.31]$): the model
\emph{looks} more monitorable without measurable faithfulness change. To probe a
substrate with dynamic range, we move to a reasoning student.

\textbf{Reasoning student: W2SR significantly degrades faithfulness}
(Table~\ref{tab:r1}). Fully in-family: student \texttt{R1-Distill-Qwen-7B}, weak
teacher \texttt{R1-Distill-Qwen-1.5B}; distillation is clean (gate format-valid
$0.99$); baseline acknowledgment is off-floor ($25\%$). W2SR drops pooled
acknowledgment to $3.2\%$ (paired on (qid,cue): $\Delta\!=\!-0.220$, CI
$[-0.293,-0.153]$, McNemar $p\!\approx\!2\times10^{-9}$, discordant $34/1$), at
\emph{preserved} GPQA accuracy ($0.42$ vs.\ $0.28$, $n{=}40$) and ${\sim}14\times$
shorter CoT. The drop is not an artifact (hand-checked: completions coherent, $37/40$
conclude with an answer, none degenerate---the model solves the problem without
verbalizing the cue).

\begin{table}[t]\centering\small
\caption{Reasoning-student monitorability (GPQA, in-family R1-distill teachers).
Faithfulness = acknowledgment among cued samples. W2SR degrades faithfulness at
preserved accuracy; a strong (14B) teacher degrades it nearly as much as a weak
(1.5B) one (compression-general).}
\label{tab:r1}
\begin{tabular}{lccc}
\toprule
condition & faithfulness & GPQA acc & median CoT (chars) \\
\midrule
baseline R1-7B & $40/160 = 25.0\%$ & 0.28 & 20{,}096 \\
W2SR, weak teacher (1.5B) & $6/190 = 3.2\%$ & 0.42 & 1{,}392 \\
W2SR, strong teacher (14B) & $13/175 = 7.4\%$ & --- & 1{,}638 \\
\bottomrule
\end{tabular}
\\[2pt]
\footnotesize weak vs.\ baseline $p\!=\!2\times10^{-9}$; strong vs.\ baseline
$p\!=\!1\times10^{-5}$; weak vs.\ strong $p\!=\!0.07$ (n.s.).
\end{table}

\textbf{Mechanism: compression-dominant + a real matched-length residual.}
Faithfulness and CoT length both fall; the near-trivial half is that shorter CoT
surfaces cues less, so the claim of interest is degradation \emph{at matched
length}. Length is endogenous to reasoning (a brevity-prompt control failed---R1
ignores it), so we isolate the residual by \emph{convergence} across imperfect
length-controlled cuts on long ($\ge\!9.6$k-char) traces: (a) acknowledgment
(judge): $19\%$ vs.\ $44\%$, $p\!=\!0.016$; (b) literal cue-mention rate: $16\%$
vs.\ $28\%$ (same direction, n.s.); (c) mention density: $0.037$ vs.\ $0.167$ per
1k chars; (d) natural length-matched question pairs: $1/9$ vs.\ $4/9$ (same
direction, underpowered). \textbf{All four cuts agree in direction; the primary
(judge) cut is significant.} So most of the $25\%\!\to\!3\%$ drop is
\textbf{compression} (W2SR's long traces recover to $19\%$), with a \textbf{real,
moderate residual} ($\sim$2.3$\times$) at matched length: the distilled model
surfaces decision-relevant cues less even per unit reasoning. We claim the residual
\emph{persists across length-controlled cuts}, not that the confound is eliminated.
A teacher-strength axis (Table~\ref{tab:r1}) shows the effect is
\textbf{compression-general}: a strong 14B teacher degrades faithfulness about as
much as the 1.5B teacher and is not distinguishable from it---W2SR (weak$\to$strong)
is not uniquely harmful; distillation-induced compression is.

\section{Discussion}
W2SR does not preserve monitorability and, where it is measurable, it hurts. The
two substrates are consistent: on the floored instruct student the composite
``monitorability'' (mean of verbosity and faithfulness) rises after distillation
almost entirely via verbosity---the model \emph{looks} more monitorable; on the
reasoning student, faithfulness has room to move and W2SR \emph{reduces} it, mostly
by compressing the CoT into a short, direct form that no longer surfaces the cue.
The safety-relevant component thus gets \emph{worse}, not better. The warning is
sharp: \emph{an accuracy-only certification of a weak-to-strong student---exactly
W2SR's---can pass a model whose chain-of-thought has become markedly less
revealing}, and a monitorability check that rewards elaborate-looking CoT would
miss it too. That the in-family requirement and headroom probe are needed for the
accuracy claim to even hold~\cite{repshape2025} only sharpens how narrow the
certified guarantee is.

\section{Limitations}
\textbf{Floored instruct substrate:} the instruct arm only motivates; the
faithfulness claim rests on the reasoning student. \textbf{Moderate / underpowered
residual:} the matched-length residual is significant on the judge measure
($p\!=\!0.016$) but directional-only on three further cuts; we claim a real,
moderate effect that persists across cuts, not an eliminated confound.
\textbf{Single family:} the degradation is shown within the Qwen R1-distill
lineage; a matched-recipe cross-family (Llama) test was confounded by
over-compression (a terse instruct teacher cratered the reasoning student's
capability), and no clean capability-matched in-family Llama teacher exists, so
cross-family generalization remains open. \textbf{Also:} LoRA (not full SFT), a
single judge model, and GPQA-Diamond only for the extension.

\section{Conclusion}
W2SR's accuracy gain is real but conditional (in-family teacher; genuine only on a
base student), and on a reasoning student it \emph{significantly reduces} how often
the CoT reveals a decisive hidden cue ($25\%\!\to\!3\%$) at preserved accuracy---
predominantly by compressing the reasoning, with a real residual beyond length, and
regardless of teacher strength. Certifying weak-to-strong students on accuracy, or
on how elaborate their CoT looks, is insufficient and potentially misleading; the
cue-acknowledgment component of monitorability must be measured directly.

\begin{thebibliography}{9}\small
\bibitem{yuan2026} Yuan et al. Incentivizing Strong Reasoning from Weak
Supervision (W2SR). arXiv:2505.20072, 2025.
\bibitem{meek2025} Meek et al. Measuring Chain-of-Thought Monitorability.
arXiv:2510.27378, 2025.
\bibitem{turpin2023} Turpin et al. Language Models Don't Always Say What They Think.
arXiv:2305.04388, 2023.
\bibitem{chua2025} Chua \& Evans. Are DeepSeek-R1 reasoning models more faithful?
arXiv:2501.08156, 2025.
\bibitem{repshape2025} Representations Shape Weak-to-Strong Generalization.
2025.
\end{thebibliography}

\end{document}
